% !PRINTED 156
\weberpage{156}

\begin{center}
\S.\,48.

\gesperrt{Resultants.}
\end{center}

Another application of the theory of symmetric functions, which moreover contains the formation of the discriminant as a special case, is the formation of resultants.

Let

\weq{1}{
f(x) = a_0\,x^n + a_1\,x^{n-1} + a_2\,x^{n-2} + \cdots + a_n,
}

\weq{2}{
\varphi(x) = b_0\,x^m + b_1\,x^{m-1} + b_2\,x^{m-2} + \cdots + b_m
}

be two integral rational functions of degree $n$ and $m$, and let $\alpha_1, \alpha_2 \ldots \alpha_n$ be the roots of the first, so that the following factorization holds:

\weq{3}{
f(x) = a_0\,(x - \alpha_1)\,(x - \alpha_2) \ldots (x - \alpha_n)
}

Now the product

\[
\varphi(\alpha_1)\,\varphi(\alpha_2) \ldots \varphi(\alpha_n)
\]

is a symmetric function of the roots of $f(x)$, and can therefore be expressed integrally and rationally in terms of the following quotients:

\[
\frac{a_1}{a_0}, \quad \frac{a_2}{a_0} \ldots \frac{a_n}{a_0}
\]

The order of this symmetric function is $(m, m \ldots m)$ (\S.\,45), and hence

\weq{4}{
R = a_0^m\,\varphi(\alpha_1)\,\varphi(\alpha_2) \ldots \varphi(\alpha_n)
}

is an integral rational and homogeneous function of order $m$ in $a_0, a_1 \ldots a_n$. It is also, as its expression shows, an integral homogeneous function of order $n$ in $b_0, b_1 \ldots b_m$. It is called the \gesperrt{resultant} of the two functions $f(x)$ and $\varphi(x)$. This designation is justified by a second representation.

If we denote the roots of $\varphi(x)$ by $\beta_1, \beta_2 \ldots \beta_m$, and thus put

\weq{5}{
\varphi(x) = b_0\,(x - \beta_1)\,(x - \beta_2) \ldots (x - \beta_m),
}

then by (4) we obtain

\weq{6}{
R = a_0^m\,b_0^n \overset{i,\,k}{\Pi}(\alpha_i - \beta_k),
}

in which the product sign $\Pi$ extends over all the indices $i = 1, 2 \ldots n$, $k = 1, 2 \ldots m$,
%CONT>>
